\documentclass[11pt]{article}
\usepackage[margin=0.9in]{geometry}
\usepackage{amsmath}
\pagestyle{empty}

\begin{document}
\section*{Canting calculation: changes}
The governing equations and material parameters are unchanged.

\begin{itemize}
\setlength{\itemsep}{0.6em}

\item Mesh: create each \verb|Gmsh| mesh once, then use custom
\verb|radial mesh motion| with fixed connectivity. This removes repeated
remeshing and the resulting jumps in computed gradients.

\item Field solves: recover continuous gradients using \verb|L2 projection|;
interpolate boundary pressure using a \verb|periodic cubic spline|.
Use \verb|LU factorization| for the linear systems and reuse the FEM setup.

\item Resolution: increase contact-line nodes from 20 to 80, use cubic FEM
fields, and decrease mesh size from $0.60$ to $0.30\,\mathrm{mm}$.
Keep \verb|BDF| time integration, with tolerances
$(\mathrm{rtol},\mathrm{atol})=(10^{-6},10^{-8})$; enforce left/right symmetry.

\item Curvature: use \verb|centered finite differences| along the contact line
$\mathbf r(s,t)$, avoiding second time derivatives of trajectory splines:
\begin{gather*}
  D_s f_i=\frac{f_{i+1}-f_{i-1}}{2\Delta s},\qquad
  \mathbf T=D_s\mathbf r,\qquad
  \mathbf v=\sigma(\theta)\mathbf n,\\
  \alpha=K_\alpha\kappa,\qquad
  \kappa=\frac{\mathbf T\times D_s\mathbf v}
                    {|\mathbf T|^2\,|\sigma(\theta)|}.
\end{gather*}
Here $\times$ denotes the scalar two-dimensional cross product.
Exclude nearly stationary samples, where curvature is unreliable.

\item Averaging: use a \verb|Gaussian kernel| and the \verb|trapezoidal rule|
for a spatial mean, replacing the time-sample-weighted robust average.
For mirror-paired samples $(+,-)$, with $\alpha$ in radians and $x$ in mm,
\begin{gather*}
  g_k=\frac{180}{\pi}\,
       \frac{\alpha_{k,+}-\alpha_{k,-}}{x_{k,+}-x_{k,-}},\qquad
  g=\frac{\sum_k W_k g_k}{\sum_k W_k},\\
  W_k=\exp\!\left[-\frac{|\mathbf r_k-\mathbf r_0|^2}{2\ell^2}\right]
      |\sigma_k|\,|\mathbf T_k|\,\omega_k,
  \qquad \ell=0.30\,\mathrm{mm}.
\end{gather*}
Here $\omega_k$ is the trapezoidal time weight; the constant material-node
spacing cancels. The evaluation point remains
$\mathbf r_0=(1.35,\,2.7\cos(\pi/6))\,\mathrm{mm}$.
Post-processing-only error bars are
$\Delta g=(\max g_{\rm variants}-\min g_{\rm variants})/2$,
varying kernel width, stationary-speed cutoff, and saved-sample spacing.

\item More runs: increase from 9 to 24 independently simulated heights,
concentrated around the bend. Use \verb|process parallelism| with six workers
and one numerical-library thread per worker. Reuse matching completed runs,
resume checkpoints, and skip completed cases. The black curve is a separate
\verb|nonlinear least squares| fit; no smoothing across heights is applied
to the simulation points.

\end{itemize}
\end{document}
